\begin{frame}{핵심 구분: ``일괄 예약'' vs ``매 스텝 재계산''}
  \begin{columns}[T]
    \begin{column}{0.48\textwidth}
      \kb{낙관적 초기화의 ``탐험''}
      \begin{itemize}
        \item \textbf{초반에 일괄 예약}(각 팔 1회).
        \item 이후 추가 탐험 장치 없음.
        \item \textbf{고정(stationary)} 환경: 최선 팔 찾기에 성공.
      \end{itemize}
    \end{column}
    \begin{column}{0.48\textwidth}
      \kb{UCB의 ``탐험''}
      \begin{itemize}
        \item \textbf{매 스텝 재계산}($N$ 이 작은 팔에 계속
          보너스 부여).
        \item \textbf{비정상 환경}: 환경 변화를 \textbf{감지하고
          적응}하는 사실상 유일한 전략.
      \end{itemize}
    \end{column}
  \end{columns}
  \vskip0.6em
  \begin{block}{Week 6(TD 학습)의 ``고정 $\alpha$ vs $1/N$'' 논쟁과
  같은 구조}
    \textbf{정보의 ``신선도''를 유지}하는지가 비정상 환경에서
    성패를 가름.
  \end{block}
\end{frame}
